% TEMPLATE-MILC-v3.tex - plantilla maestra UNICA de las guias MILC v3.
%
% La usan las tres familias de guias, cada una con su driver:
%   · Grados 6-11  -> scripts/build-guias-g11.py
%   · Bebras       -> scripts/build-guias-bebras.py
%   · Semillero    -> scripts/build-guias-semillero.py
%
% Antes habia tres copias casi identicas (~400 lineas iguales, 6 distintas):
% cada mejora habia que aplicarla tres veces, y en la practica no se hacia
% (el motor de recursos vivia solo en dos de ellas). Lo que varia entre
% familias esta parametrizado en 6 placeholders que cada driver llena
% (se nombran aqui SIN su sintaxis real: el builder reemplaza por texto plano,
% tambien dentro de los comentarios, y un valor multilinea rompe el preambulo):
%
%   HEADER_IZQ        encabezado izquierdo de pagina
%   HEADER_DER        encabezado derecho de pagina
%   PORTADA_ETIQUETA  rotulo sobre el numero grande (GRADO/MOMENTO/MODULO)
%   PORTADA_SERIE     linea de serie de la portada
%   PORTADA_ID        identificador de la guia en la portada
%   PORTADA_CIRCULO   el circulito de grado (°), o vacio si no aplica
%
% El resto de claves las documenta content/guias/_SCHEMA.md. El script valida
% que no queden placeholders sin reemplazar antes de escribir el archivo final.

\documentclass[11pt,letterpaper]{article}
\usepackage[margin=2.0cm]{geometry}
\usepackage[table]{xcolor}
\usepackage{fontspec}
\usepackage[spanish]{babel}
\usepackage{tikz}
% Librerías TikZ para los diagramas de las guías (bloque `recursos:`):
%  · babel  → babel en español vuelve activos caracteres como «>», y TikZ los usa
%             en sus opciones (>=stealth, ->). Sin esta librería, cualquier
%             diagrama con flechas revienta con «\language@active@arg> has an
%             extra }». Debe ir ANTES de que se dibuje nada.
%  · positioning        → sintaxis «below left=2pt and 3pt».
%  · decorations.pathreplacing → llaves ({brace}) para señalar rangos.
\usetikzlibrary{babel,positioning,decorations.pathreplacing}
\usepackage{graphicx}
% Circuitos: el trabajo de electrónica se hace hoy con micro:bit y sus sensores
% integrados (MakeCode, sin cableado), así que no se necesitan esquemáticos.
% Si algún día entran montajes con protoboard, basta descomentar esta línea:
% los diagramas .tex/.tikz declarados en `recursos:` ya se incrustan solos.
% \usepackage{circuitikz}
\usepackage[most]{tcolorbox}
\usepackage{tabularx}
\usepackage{array}
\usepackage{fancyhdr}
\usepackage{titlesec}
\usepackage[document]{ragged2e}  % bandera a la derecha en todo el cuerpo
\usepackage{enumitem}
\usepackage{microtype}

% Helvetica Neue no trae la flecha U+2192, asi que cada «→» escrito literalmente
% en el YAML se compilaba a NADA, en silencio (462 flechas en 61 guias: rutas del
% tipo «paleta → tipografia → lockup» salian rotas). Mapeamos el caracter a su
% version matematica, que si tiene glifo. Asi el autor puede seguir escribiendo
% «→» comodamente en el YAML.
\usepackage{newunicodechar}
\newunicodechar{→}{\ensuremath{\rightarrow}}
% Mismo problema con los simbolos de marca y vinetas: el «✓» de «una palomita ✓»
% salia como cuadrito vacio en el PDF de todas las guias que lo usaban.
\usepackage{pifont}
\newunicodechar{✓}{\ding{51}}
\newunicodechar{✗}{\ding{55}}
\newunicodechar{✔}{\ding{52}}
\newunicodechar{•}{\textbullet}
\newunicodechar{≥}{\ensuremath{\geq}}
\newunicodechar{≤}{\ensuremath{\leq}}

\setmainfont{Helvetica Neue}
\setsansfont{Helvetica Neue}
\setlength{\parindent}{0pt}
\setlength{\parskip}{6pt}
% Sin particion de palabras: en 6.º-7.º una palabra cortada por guion al final
% de linea («Comuni-cacion») frena la lectura mucho mas que una linea corta.
\hyphenpenalty=10000
\exhyphenpenalty=10000
\pretolerance=2000
\tolerance=3000
\setlength{\emergencystretch}{3em}
\linespread{1.10}
\renewcommand{\arraystretch}{1.25}
\setlist{leftmargin=5mm,itemsep=1mm,topsep=1mm}
\newcolumntype{Y}{>{\RaggedRight\arraybackslash}X}

\definecolor{milcMagenta}{HTML}{E6007E}
\definecolor{milcVino}{HTML}{5A0038}
\definecolor{milcTurquesa}{HTML}{0093A5}
\definecolor{milcVerde}{HTML}{58A923}
\definecolor{milcMostaza}{HTML}{E5B400}
\definecolor{milcOcre}{HTML}{B86600}
\definecolor{milcNegro}{HTML}{241816}
\definecolor{milcGris}{HTML}{F7F7F4}
\definecolor{milcLinea}{HTML}{E8E2DD}
\definecolor{milcGradePrimary}{HTML}{84CC16}
\definecolor{milcGradeDark}{HTML}{4D7C0F}
\definecolor{milcGradeSoft}{HTML}{ECFCCB}
\definecolor{milcGradeAccent}{HTML}{CFE7FF}
\definecolor{milcGradeText}{HTML}{FFFFFF}
\color{milcNegro}

\pagestyle{fancy}
\fancyhf{}
\lhead{\small Tecnología e Informática · Grado 7}
\rhead{\small Guía 13 · MILC v3}
\cfoot{\small\thepage}
\renewcommand{\headrulewidth}{0pt}

\newtcolorbox{titlebox}[1]{
  enhanced,colback=#1,colframe=#1,arc=7mm,boxrule=0pt,
  left=7mm,right=7mm,top=3mm,bottom=3mm,
  before skip=8pt,after skip=8pt,
  fontupper=\sffamily\bfseries\Large\color{white}
}

\newtcolorbox{softbox}[3]{
  enhanced,colback=#2,colframe=#1,borderline west={4pt}{0pt}{#1},
  arc=2mm,boxrule=.4pt,left=4mm,right=4mm,top=3mm,bottom=3mm,
  before skip=6pt,after skip=8pt,title={#3},coltitle=white,
  colbacktitle=#1,fonttitle=\sffamily\bfseries,fontupper=\RaggedRight,
  attach boxed title to top left={xshift=3mm,yshift=-2mm},
  boxed title style={arc=2mm, boxrule=0pt}
}

\newtcolorbox{drawbox}[2]{
  enhanced,colback=white,colframe=#1,arc=4mm,boxrule=1pt,
  left=5mm,right=5mm,top=4mm,bottom=4mm,before skip=8pt,after skip=8pt,
  title={#2},coltitle=white,colbacktitle=#1,fonttitle=\sffamily\bfseries,
  fontupper=\RaggedRight,
  attach boxed title to top left={xshift=4mm,yshift=-2mm},
  boxed title style={arc=3mm, boxrule=0pt}
}

\newtcolorbox{infoband}[1]{
  enhanced,colback=#1!8,colframe=#1!50,arc=2mm,boxrule=.5pt,
  left=4mm,right=4mm,top=2mm,bottom=2mm,before skip=4pt,after skip=4pt,
  fontupper=\small\RaggedRight
}

% Puente narrativo entre secciones (contrato editorial Conéctate).
% Da continuidad: "Acabas de X. Ahora vas a Y porque Z."
% Visualmente: banda izquierda en color del grado, texto en itálica pequeña.
\newtcolorbox{puentebox}{
  enhanced,colback=milcGradeSoft,colframe=milcGradeDark,
  borderline west={3pt}{0pt}{milcGradeDark},
  arc=0mm,boxrule=0pt,left=5mm,right=4mm,top=2.5mm,bottom=2.5mm,
  before skip=4pt,after skip=8pt,
  fontupper=\itshape\small\color{milcNegro}\RaggedRight
}
% La flecha va en modo matematico a proposito: Helvetica Neue NO tiene el glifo
% U+2192, asi que \textrightarrow se compilaba a NADA («Missing character: There
% is no → in font Helvetica Neue») y el puente salia sin flecha. En modo
% matematico la flecha viene de la fuente matematica y si se dibuja.
\newcommand{\puente}[1]{%
  \begin{puentebox}%
    \textcolor{milcGradeDark}{$\rightarrow$}\hspace{2mm}#1%
  \end{puentebox}%
}

\newcommand{\linead}[1]{\noindent\color{milcLinea}\rule{#1}{0.5pt}\color{milcNegro}}
\newcommand{\respuesta}[1]{\par\vspace{2mm}\foreach \n in {1,...,#1}{\noindent\color{milcLinea}\rule{\linewidth}{0.5pt}\par\vspace{5mm}}\color{milcNegro}}
\newcommand{\checkbox}{\textcolor{milcGradeDark}{\fbox{\phantom{x}}}\hspace{5pt}}

\newcommand{\phasecard}[4]{
\begin{tcolorbox}[enhanced,colback=#3,colframe=#2,arc=3mm,boxrule=.5pt,height=3.05cm,valign=center,halign=center,left=1.5mm,right=1.5mm,top=2mm,bottom=2mm]
{\sffamily\bfseries\fontsize{22}{24}\selectfont\textcolor{#2}{#1}}\par
{\sffamily\bfseries\small #4}
\end{tcolorbox}
}

% ── Piezas del rediseno 2026 (legibilidad 6.º-11.º) ───────────────────
% Banda de estacion: reemplaza el titlebox macizo. Titulo a la izquierda,
% dato practico (tiempo/modalidad) a la derecha, en una sola linea.
\newcommand{\milcestacion}[2]{%
\begin{tcolorbox}[enhanced,colback=#1,colframe=#1,arc=2mm,boxrule=0pt,
  left=5mm,right=5mm,top=2.6mm,bottom=2.6mm,before skip=11pt,after skip=7pt]
{\sffamily\bfseries\fontsize{15}{18}\selectfont\color{white}#2}
\end{tcolorbox}}

% Tarjeta de paso numerada: sustituye las tablas «Pilar / Aplicacion», que
% eran formato de consulta docente, no de estudiante. El numero va en un
% circulo sobre el borde para que la secuencia se lea de un vistazo.
\newcommand{\milcpaso}[3]{%
\begin{tcolorbox}[enhanced,colback=white,colframe=milcLinea,arc=1.5mm,boxrule=.9pt,
  left=14mm,right=4mm,top=3mm,bottom=3mm,before skip=4pt,after skip=4pt,
  fontupper=\RaggedRight,
  overlay={\node[anchor=north west,circle,fill=#2,text=white,inner sep=0pt,
    minimum size=7.4mm,font=\sffamily\bfseries\small]
    at ([xshift=3.6mm,yshift=-3.2mm]frame.north west) {#1};}]
#3
\end{tcolorbox}}

% Casilla marcable de tamano util con lapiz (la anterior era \fbox{\phantom{x}},
% de alto variable segun el texto que la seguia).
\newcommand{\milcmarca}{\framebox[3.8mm]{\rule{0pt}{3.4mm}}}

% Fila marcable de lista suelta (checks de la estacion 1).
\newcommand{\milccheckrow}[1]{%
\par\vspace{1.8mm}\noindent
\makebox[6.4mm][l]{\raisebox{-.55mm}{\textcolor{milcNegro}{\milcmarca}}}%
\parbox[t]{\dimexpr\linewidth-6.4mm\relax}{\RaggedRight #1}\par}

% Celda marcable dentro de la rubrica (tabla de autoevaluacion). Misma
% sangria francesa, pero pensada para vivir dentro de una columna X.
\newcommand{\milcopcion}[1]{%
\makebox[5.2mm][l]{\raisebox{-.55mm}{\textcolor{milcNegro}{\milcmarca}}}%
\parbox[t]{\dimexpr\linewidth-5.2mm\relax}{\RaggedRight #1}}

% ── Recursos visuales de la guía ──────────────────────────────────────
% Declarados en el bloque `recursos:` del YAML y emitidos por el builder.
% \guiaFigura   → imagen rasterizada o PDF (foto, carta celeste, montaje).
% \guiaDiagrama → fuente TikZ/circuitikz (.tex) incrustada como vector:
%                 es la vía para circuitos y esquemáticos editables.
% #1 = ruta relativa al .tex · #2 = pie (puede ir vacío)
\newcommand{\guiaFigura}[2]{%
\begin{center}
\includegraphics[width=0.88\linewidth,height=0.38\textheight,keepaspectratio]{#1}\\[1.5mm]
{\footnotesize\itshape\textcolor{milcNegro!75}{#2}}
\end{center}
}

\newcommand{\guiaDiagrama}[2]{%
\begin{center}
\input{#1}\\[1.5mm]
{\footnotesize\itshape\textcolor{milcNegro!75}{#2}}
\end{center}
}

\begin{document}
\thispagestyle{empty}

% ── Portada ───────────────────────────────────────────────────────────
% Antes: un rectangulo de color a pagina completa. Se veia bien en pantalla,
% pero en la fotocopiadora del colegio se comia el toner y salia gris sucio.
% Ahora la banda ocupa el tercio superior y el resto va en blanco.
\begin{tikzpicture}[remember picture,overlay]
  \path[fill=milcGradePrimary]
    (current page.north west) rectangle ([yshift=-10.6cm]current page.north east);

  \node[anchor=north west,text=milcGradeText] at ([xshift=2.0cm,yshift=-1.55cm]current page.north west)
    {\sffamily\bfseries\fontsize{12}{14}\selectfont GRADO};
  \node[anchor=north west,text=milcGradeText,opacity=.85] at ([xshift=2.0cm,yshift=-2.35cm]current page.north west)
    {\sffamily\bfseries\fontsize{8.5}{10}\selectfont SERIE GUÍAS · TECNOLOGÍA E INFORMÁTICA};
  \node[anchor=north east,text=milcGradeText,opacity=.85,align=right] at ([xshift=-2.0cm,yshift=-1.55cm]current page.north east)
    {\sffamily\bfseries\fontsize{9}{12}\selectfont I.E. Sor María Juliana\\Cartago, Valle del Cauca};

  \node[anchor=west,text=milcGradeText] (gradonum) at ([xshift=1.85cm,yshift=-7.1cm]current page.north west)
    {\sffamily\bfseries\fontsize{112}{112}\selectfont 7};
  % El circulito de grado (°) solo aplica a las guias de grado («11°»); en
  % Bebras y Semillero el numero grande es un momento/modulo, no un grado.
  % Se posiciona relativo al nodo `gradonum`, no en coordenadas absolutas.
  \draw[milcGradeText,line width=1.6pt]
    ([xshift=2.0mm,yshift=-4.5mm]gradonum.north east) circle (.15cm);

  \node[anchor=west,text=milcGradeText,text width=11.4cm,align=left] at ([xshift=1.5cm,yshift=0.5cm]gradonum.east)
    {
     {\sffamily\bfseries\fontsize{9}{11}\selectfont\MakeUppercase{Período 2 · Algoritmia y pensamiento computacional}}\\[3mm]
     {\sffamily\bfseries\fontsize{21}{25}\selectfont\setlength{\baselineskip}{27pt}Pensamiento\\[4pt]computacional ---\\[4pt]las 4 técnicas\\[4pt]del oficio}\\[3.5mm]
     {\sffamily\bfseries\fontsize{15}{18}\selectfont Guía 13-7-TIC}
    };
\end{tikzpicture}

\vspace*{9.4cm}
\vfill

{\sffamily\bfseries\fontsize{8}{10}\selectfont\textcolor{milcGradeDark}{LO QUE TE LLEVAS HOY}}

\begin{tcolorbox}[enhanced,colback=white,colframe=milcNegro,arc=2mm,boxrule=1.6pt,
  left=5mm,right=5mm,top=4mm,bottom=4mm,before skip=3pt,after skip=12pt,
  fontupper=\RaggedRight\fontsize{11}{15.5}\selectfont]
Una \textbf{aplicación de las 4 técnicas} a un problema real escogido por ti (organizar una fiesta de cumpleaños, planear un viaje al pueblo, hacer un trabajo de Sociales, etc.). En el cuaderno: \textbf{(a) descomposición} del problema en sub-problemas (mínimo 4); \textbf{(b) patrones} identificados; \textbf{(c) abstracción} aplicada (qué dejas afuera porque no es esencial); \textbf{(d) algoritmo final} en pasos numerados. Más \textbf{una reflexión} sobre cuál técnica te costó más.
\end{tcolorbox}

\begin{tcolorbox}[enhanced,colback=milcGris,colframe=milcGris,arc=2mm,boxrule=0pt,
  left=5mm,right=5mm,top=3.5mm,bottom=3.5mm,after skip=12pt]
{\sffamily\bfseries\fontsize{8}{10}\selectfont\textcolor{milcNegro!55}{ESTA GUÍA ES DE}}\\[3mm]
\begin{tabularx}{\linewidth}{X p{3.2cm} p{3.2cm}}
\linead{\linewidth} & \linead{3.0cm} & \linead{3.0cm}\\[-1mm]
{\fontsize{8.5}{10}\selectfont\textcolor{milcNegro!55}{Nombre completo}} &
{\fontsize{8.5}{10}\selectfont\textcolor{milcNegro!55}{Grupo}} &
{\fontsize{8.5}{10}\selectfont\textcolor{milcNegro!55}{Fecha}}\\
\end{tabularx}
\end{tcolorbox}

\vfill

\noindent\textcolor{milcLinea}{\rule{\linewidth}{1pt}}\\[2mm]
{\sffamily\bfseries\fontsize{9.5}{12}\selectfont Autor: Dr. Álvaro Cárdenas Orozco}\\[1mm]
{\sffamily\fontsize{8.5}{11}\selectfont\textcolor{milcNegro!55}{Metodología MILC · apertura ancestral · pensamiento computacional · triángulo Dussel–estoicismo–Floridi}}

\newpage

\section*{Apertura: Pensamiento computacional --- las 4 técnicas del oficio}

\begin{softbox}{milcGradeDark}{milcGradeSoft}{Saber ancestral}
\textbf{Cuando una abuela campesina del Valle del Cauca preparaba el sancocho para 30 personas, no se ponía a llorar ante la inmensidad de la tarea: la dividía en partes.} \emph{``Una persona pela papas''}, decía. \emph{``Otra pela yuca''}. \emph{``Otra lava el plátano''}. \emph{``Otra prende el fogón''}. \emph{``Otra trae el agua del nacedero''}. En 1 hora, el sancocho está listo. \textbf{Esa habilidad --- dividir un problema grande en sub-problemas manejables --- es la primera técnica del pensamiento computacional}. Se llama \textbf{descomposición}. Pero la abuela también notaba algo más: que \emph{pelar papa y pelar yuca son tareas parecidas} (las dos requieren cuchillo, paciencia, retirar la cáscara). Eso es \textbf{patrones}. Y notaba que \emph{no importaba si la papa tenía manchas o no: la cocinaba igual} --- dejaba afuera los detalles irrelevantes. Eso es \textbf{abstracción}. Y al final, tenía \emph{el orden exacto de pasos} para que todo saliera al mismo tiempo. Eso es \textbf{algoritmo}. \emph{Las 4 técnicas del pensamiento computacional las usaban las abuelas para hacer sancocho, mucho antes de que existieran los computadores}.
\end{softbox}

\begin{softbox}{milcTurquesa}{milcGris}{Saber contemporáneo}
El \textbf{pensamiento computacional} no es exclusivo de la programación: es una \emph{forma de abordar problemas} que se aplica en cualquier campo. Su definición moderna viene de \textbf{Jeannette Wing} (científica computacional de Carnegie Mellon, 2006). Tiene \textbf{4 técnicas}: \textbf{(1) Descomposición}: dividir un problema grande en sub-problemas pequeños. \textbf{(2) Reconocimiento de patrones}: identificar similitudes entre sub-problemas y reusar soluciones. \textbf{(3) Abstracción}: dejar afuera los detalles que no son esenciales para enfocarse en lo importante. \textbf{(4) Diseño de algoritmos}: armar la secuencia de pasos que resuelve el problema. \textbf{Las 4 funcionan juntas:} descompones para entender, reconoces patrones para no repetir trabajo, abstraes para simplificar, diseñas el algoritmo final. \emph{Aplicar las 4 técnicas a un problema cualquiera te entrena en una forma de pensar que vale para programar, estudiar, trabajar, vivir}.
\end{softbox}

\begin{softbox}{milcMostaza}{milcGris}{Pregunta-puente}
\textit{Si te pidieran organizar la \emph{fiesta de cumpleaños número 15} de tu prima (50 personas, presupuesto limitado, 3 días para organizarla), ¿\textbf{por dónde empezarías}? ¿Te paralizarías o sabrías descomponer el problema?}
\end{softbox}

\begin{softbox}{milcVerde}{milcGris}{Saber y saber hacer}
Al terminar la clase: (1) podrás \textbf{identificar} las 4 técnicas del pensamiento computacional; (2) sabrás \textbf{aplicar} cada una a un problema real; (3) podrás \textbf{evaluar} qué técnica necesita un problema en cada momento; (4) habrás \textbf{creado} una aplicación completa de las 4 técnicas a un problema escogido por ti.
\end{softbox}



\small
\begin{tabularx}{\textwidth}{>{\bfseries}p{3.0cm}Y}
\rowcolor{milcGradePrimary!12}Autor & Dr. Álvaro Cárdenas Orozco\\
DBA & Construye algoritmos y diagramas de flujo para resolver problemas paso a paso (MEN, T\&I 7°)\\
Referentes & Saber ancestral del tejedor · Pensamiento computacional · Scratch\\
Producto final & Una \textbf{aplicación de las 4 técnicas} a un problema real escogido por ti (organizar una fiesta de cumpleaños, planear un viaje al pueblo, hacer un trabajo de Sociales, etc.). En el cuaderno: \textbf{(a) descomposición} del problema en sub-problemas (mínimo 4); \textbf{(b) patrones} identificados; \textbf{(c) abstracción} aplicada (qué dejas afuera porque no es esencial); \textbf{(d) algoritmo final} en pasos numerados. Más \textbf{una reflexión} sobre cuál técnica te costó más.\\
\end{tabularx}
\normalsize

\puente{Hoy haces 4 cosas: aprendes las 4 técnicas con ejemplos, escoges un problema real, aplicas las 4 técnicas, evalúas el proceso.}

\milcestacion{milcGradeDark}{Ruta MILC: así vamos a aprender}
\begin{tabularx}{\textwidth}{>{\centering\arraybackslash}X>{\centering\arraybackslash}X>{\centering\arraybackslash}X>{\centering\arraybackslash}X}
\phasecard{1}{milcVerde}{milcGris}{Escucha\\[-1mm]\small Observo, escucho y pregunto.} &
\phasecard{2}{milcTurquesa}{milcGris}{Entender\\[-1mm]\small Sistematizo y ordeno.} &
\phasecard{3}{milcMagenta}{milcGris}{Hacer\\[-1mm]\small Construyo y aplico.} &
\phasecard{4}{milcVino}{milcGris}{Cierre\\[-1mm]\small Evalúo y reflexiono.}
\end{tabularx}


\puente{Empezamos con un problema grande para ver por qué da pavor sin descomponer.}

\milcestacion{milcVerde}{Estación 1 de 3 · Escucha}

\begin{softbox}{milcVerde}{milcGris}{Escena para escuchar}
\textbf{Actividad 1 · IDENTIFICA --- Un problema grande para pensar} (10 min · individual). Lee este caso y en el cuaderno responde solo \textbf{cómo te sentirías} si te lo asignaran. \textbf{El caso:} \emph{``Eres responsable de organizar la celebración de los 50 años de tu colegio. Tienes 60 días, 800 estudiantes, presupuesto limitado, debes coordinar con padres, profesores, exalumnos, prensa local, alcaldía. Hay que hacer: ceremonia, decoración, sonido, comida, invitados, prensa, recordatorios, fotos.''}. \emph{En el cuaderno responde:} \emph{(1) ¿Cómo te sientes ante este problema? (abrumado, motivado, paralizado, etc.)}. \emph{(2) ¿Por dónde empezarías? (cuenta cómo)}. \emph{(3) ¿Crees que se puede hacer solo, o necesitas equipo?}. \emph{(4) ¿Cuánto tiempo crees que tarda planearlo bien?}.
\end{softbox}

{\sffamily\bfseries\fontsize{8}{10}\selectfont\textcolor{milcNegro!55}{MARCA CUANDO LO HAYAS HECHO}}
\milccheckrow{Leí el caso completo.}
\milccheckrow{Reconocí mis emociones honestas ante un problema grande.}
\milccheckrow{Respondí las 4 preguntas con sinceridad.}

\begin{infoband}{milcVerde}


\textbf{Lo que va al cuaderno (Actividad 1 · IDENTIFICA).} Título: «Actividad 1 --- Reacción ante problema grande». \textbf{Formato:} 4 respuestas. \textbf{Extensión:} 5-6 líneas. \textbf{Sabes que terminaste cuando} reconoces que sin método, el problema da pavor.
\end{infoband}

\puente{Después aprendes las 4 técnicas con sus ejemplos.}

\milcestacion{milcTurquesa}{Estación 2 de 3 · Entender}

\begin{softbox}{milcTurquesa}{milcGris}{Lo que vas a entender}
Tu reacción honesta es válida: \textbf{los problemas grandes paralizan} si no se sabe descomponerlos. El pensamiento computacional te da el método. Aprende las 4 técnicas con ejemplos del caso del colegio.
\end{softbox}

\milcpaso{1}{milcTurquesa}{\textbf{Técnica 1 --- DESCOMPOSICIÓN.} Dividir el problema grande en sub-problemas. \emph{``Celebrar los 50 años''} se divide en: \emph{(a) Ceremonia formal}, \emph{(b) Decoración del colegio}, \emph{(c) Sonido y tarima}, \emph{(d) Comida y bebida}, \emph{(e) Invitar a exalumnos y prensa}, \emph{(f) Souvenirs y recordatorios}, \emph{(g) Plan B en caso de lluvia}. \textbf{Cada sub-problema es manejable; el problema grande no lo era}. Esta es la técnica que más alivia la sensación de abrumamiento.}
\milcpaso{2}{milcTurquesa}{\textbf{Técnica 2 --- PATRONES.} Identificar similitudes entre sub-problemas para no inventar de cero cada vez. \emph{``Decoración del colegio''} y \emph{``Decoración de las aulas''} son problemas parecidos: se puede usar el mismo equipo, las mismas reglas. \emph{``Invitar a exalumnos''} y \emph{``Invitar a prensa''} son parecidos: ambos requieren lista, correos, confirmaciones. \textbf{Reconocer patrones te ahorra tiempo}: si resolviste uno, resuelves el otro. Las abuelas que pelaban papa y yuca al tiempo aplicaban patrones.}
\milcpaso{3}{milcTurquesa}{\textbf{Técnica 3 --- ABSTRACCIÓN.} Dejar afuera los detalles que no son esenciales en este momento. Cuando planeas la ceremonia, no te importa todavía el menú exacto del almuerzo. Cuando armas la lista de invitados, no te importa todavía el color de los manteles. \textbf{Abstraer significa enfocarse en lo importante de cada sub-problema, sin perderse en detalles que vendrán después}. Sin abstracción te abrumas; con ella, avanzas.}
\milcpaso{4}{milcTurquesa}{\textbf{Técnica 4 --- ALGORITMO.} Después de descomponer, encontrar patrones y abstraer, armas la \emph{secuencia exacta de pasos} para resolver. \emph{``Día 1-15: planeación general. Día 16-30: invitaciones y confirmaciones. Día 31-45: producción de decoración y comida. Día 46-55: ensayo y prueba. Día 56-60: ejecución''}. Cada sub-problema tiene su mini-algoritmo dentro del algoritmo grande. Esta es la culminación: el método operativo que se puede ejecutar.}

\begin{drawbox}{milcTurquesa}{4 técnicas + cómo funcionan juntas}
\textbf{1.} DESCOMPOSICIÓN --- dividir en sub-problemas. \\[1mm]
\textbf{2.} PATRONES --- reconocer similitudes. \\[1mm]
\textbf{3.} ABSTRACCIÓN --- dejar lo no esencial afuera. \\[1mm]
\textbf{4.} ALGORITMO --- secuencia de pasos final. \\[1mm]
\textbf{Funcionan juntas:} descompones, encuentras patrones, abstraes, armas algoritmo.
\end{drawbox}

\begin{softbox}{milcMostaza}{milcGris}{Errores comunes y su corrección}
\textbf{Errores frecuentes al aplicar las 4 técnicas.} (1) \emph{Saltarse la descomposición} --- intentar resolver el problema completo. Se vuelve abrumador. (2) \emph{No buscar patrones} --- inventar cada sub-problema desde cero, perdiendo tiempo. (3) \emph{No abstraer} --- preocuparse por detalles irrelevantes al inicio. La ceremonia primero, los manteles después. (4) \emph{Hacer algoritmo sin descomponer} --- el algoritmo queda inmanejable. (5) \emph{Creer que las técnicas son lineales} --- en realidad se aplican en ciclos. Descompones, ves patrones, abstraes, y a veces vuelves a descomponer más fino.
\end{softbox}

\begin{infoband}{milcTurquesa}


\textbf{Lo que va al cuaderno (Actividad 2 · EXPLICA).} Título: «Actividad 2 --- 4 técnicas + ejemplo». \textbf{Formato:} bloque A con las 4 técnicas definidas + 1 ejemplo cada una. Bloque B con un ejemplo del caso del cumpleaños 15 aplicando las 4. \textbf{Extensión:} 10 líneas. \textbf{Sabes que terminaste cuando} podrías explicarle las 4 técnicas a un compañero usando un ejemplo cotidiano.
\end{infoband}

\puente{Con todo claro, aplicas las 4 técnicas a un problema tuyo.}

\milcestacion{milcMagenta}{Estación 3 de 3 · Hacer}

\begin{softbox}{milcMagenta}{milcGris}{Cómo vas a construir tu producto}
\textbf{Actividad 3 · ANALIZA --- Aplica las 4 técnicas a un problema tuyo} (32 min · individual). Vas a tomar un problema real que tengas (no académico, no inventado) y aplicar las 4 técnicas paso a paso.
\end{softbox}

\milcpaso{1}{milcMagenta}{\textbf{Paso 1 --- Escoge tu problema.} En el cuaderno escribe: \emph{``Mi problema: \_\_\_\_\_''}. Tiene que ser \emph{real} (no inventado), \emph{relativamente grande} (que se sienta abrumador), y \emph{que esté en tu radar} (algo que ya quieres hacer o que tienes pendiente). Sugerencias: \emph{organizar tu cuarto antes de que llegue la abuela}, \emph{planear una salida con amigos al parque}, \emph{preparar para un examen de varias materias}, \emph{ayudar a tu mamá con la mudanza}, \emph{lanzar un mini-emprendimiento de venta de algo en el colegio}.}
\milcpaso{2}{milcMagenta}{\textbf{Paso 2 --- Descomposición.} Subtítulo: \textbf{``Mi descomposición''}. Lista numerada de \textbf{mínimo 4 sub-problemas} en que se divide tu problema grande. Cada sub-problema debe ser manejable, no abrumador. Ejemplo: si el problema es \emph{``organizar mi cuarto antes de que llegue la abuela''}, los sub-problemas podrían ser: \emph{(1) Recoger ropa sucia y limpia; (2) Ordenar la biblioteca; (3) Barrer y trapear; (4) Limpiar el escritorio; (5) Cambiar las sábanas; (6) Sacar la basura}.}
\milcpaso{3}{milcMagenta}{\textbf{Paso 3 --- Patrones.} Subtítulo: \textbf{``Patrones que encuentro''}. Identifica \textbf{2-3 similitudes} entre tus sub-problemas. Ejemplo: \emph{Recoger ropa, ordenar biblioteca y limpiar escritorio son ``ordenar superficies''. Barrer y trapear son ``limpieza del piso''. Cambiar sábanas y sacar basura son ``actividades de cierre''}. Identificar patrones agrupa tareas que se pueden hacer con la misma estrategia.}
\milcpaso{4}{milcMagenta}{\textbf{Paso 4 --- Abstracción.} Subtítulo: \textbf{``Lo que dejo afuera por ahora''}. Lista 3-4 detalles que NO son esenciales en este momento y que decides dejar para después o ignorar. Ejemplo: \emph{``No me preocupo por la decoración perfecta; no me importa qué libro vaya antes de cuál; no necesito hacer limpieza profunda --- es solo organizar''}. La abstracción te enfoca.}
\milcpaso{5}{milcMagenta}{\textbf{Paso 5 --- Algoritmo final.} Subtítulo: \textbf{``Mi algoritmo''}. Lista numerada de \textbf{8-12 pasos concretos} para ejecutar el plan. Aplica orden lógico (en el ejemplo: \emph{primero recoger ropa, después ordenar biblioteca, luego escritorio, después cambiar sábanas, finalmente barrer y trapear y sacar basura --- en ese orden para no tener que limpiar dos veces}).}
\milcpaso{6}{milcMagenta}{\textbf{Paso 6 --- Reflexión final.} Debajo del algoritmo, escribe: \emph{``La técnica que más me costó aplicar fue: \_\_\_\_\_ porque \_\_\_\_\_''}. \emph{``Si no hubiera aplicado el pensamiento computacional, habría hecho: \_\_\_\_\_''}. Firma con nombre y fecha. La reflexión es lo que te muestra cuánto creciste hoy.}

\begin{drawbox}{milcMagenta}{Antes de cerrar la S3}
\checkbox Tengo problema real escogido (no inventado). \\[1mm]
\checkbox Descomposición con mínimo 4 sub-problemas. \\[1mm]
\checkbox Identifiqué 2-3 patrones entre sub-problemas. \\[1mm]
\checkbox Hice abstracción (3-4 detalles dejados afuera). \\[1mm]
\checkbox Algoritmo final con 8-12 pasos numerados. \\[1mm]
\checkbox Reflexión final firmada.
\end{drawbox}

\begin{softbox}{milcMagenta}{milcGris}{Plantilla de trabajo}
\textbf{Estructura.} \textit{Paso 1:} mi problema. \textit{Paso 2:} descomposición (4+ sub-problemas). \textit{Paso 3:} patrones (2-3 similitudes). \textit{Paso 4:} abstracción (3-4 detalles dejados afuera). \textit{Paso 5:} algoritmo (8-12 pasos). \textit{Paso 6:} reflexión + firma.
\end{softbox}

\begin{infoband}{milcMagenta}


\textbf{Lo que va al cuaderno (Actividad 3 · ANALIZA).} Título: «Actividad 3 --- 4 técnicas aplicadas a [mi problema]». \textbf{Formato:} 5 secciones (problema + descomposición + patrones + abstracción + algoritmo) + reflexión. \textbf{Extensión:} 1 página completa. \textbf{Sabes que terminaste cuando} sientes que el problema dejó de abrumarte porque tienes método.
\end{infoband}

\puente{El producto es esa aplicación + reflexión final.}

\milcestacion{milcMostaza}{Producto de la sesión}
\textbf{4 técnicas aplicadas a problema real · descomposición + patrones + abstracción + algoritmo · firmado}.
\begin{softbox}{milcMostaza}{milcGris}{Criterios de calidad}
(1) El problema escogido es real (no inventado). (2) Descomposición tiene mínimo 4 sub-problemas manejables. (3) Hay 2-3 patrones identificados. (4) Hay 3-4 detalles abstraídos (dejados afuera). (5) El algoritmo final tiene 8-12 pasos numerados.
\end{softbox}

\puente{Antes de salir, verifica que aplicaste las 4 técnicas, no solo 1 o 2.}

\milcestacion{milcVino}{Evaluación liberadora · cinco dimensiones}

\begin{softbox}{milcVino}{milcGris}{Sentido de la evaluación}
Evaluar no es solo revisar una nota. Es reconocer qué aprendí, qué cambió en mí, qué emociones manejé, cómo me relaciono con la ciudad y la comunidad, y qué le dejaré a quien viene detrás.
\end{softbox}

\textbf{Mi reflexión liberadora en cinco dimensiones.} Completa cada frase iniciada:

\small
\begin{tabularx}{\textwidth}{>{\bfseries\RaggedRight\arraybackslash}p{3.4cm}Y>{\RaggedRight\arraybackslash}p{4.6cm}}
\rowcolor{milcVino}\color{white}Dimensión & \color{white}\bfseries Mi respuesta (frame starter) & \color{white}\bfseries Algo que descubrí\\
1. Personal & Lo que más cambió en mí fue \linead{6.5cm} & \linead{4.2cm}\\
2. Emocional & La emoción más fuerte fue \linead{2.5cm} y la regulé \linead{3cm} & \linead{4.2cm}\\
3. Ciudadana & En democracia esto importa porque \linead{6cm} & \linead{4.2cm}\\
4. Local (Cartago) & En mi barrio sería útil para \linead{6.5cm} & \linead{4.2cm}\\
5. Intergeneracional & A mi abuela/abuelo le diría \linead{6.5cm} & \linead{4.2cm}\\
\end{tabularx}
\normalsize

\begin{infoband}{milcVino}
\textbf{Para la próxima sustentación.} Vuelve a esta tabla en seis meses. Verás que las respuestas cambian — no porque lo aprendido cambie, sino porque tú cambias. La evaluación liberadora es una conversación contigo a través del tiempo.
\end{infoband}


\puente{Cerramos con 3 voces sobre por qué pensar en 4 técnicas te hace más capaz frente a problemas grandes.}

\milcestacion{milcGradeDark}{Triángulo de pensamiento · cierre filosófico}

\begin{softbox}{milcVino}{milcGris}{Dussel · lente del nosotros}
\textit{Lo que parece imposible solo, se vuelve posible cuando lo divides en partes. Ese saber salva vidas.}

Las personas que enfrentan tragedias --- una mudanza forzada, la pérdida de un trabajo, una enfermedad de un familiar --- a veces \textbf{se paralizan} ante la inmensidad. La técnica de \emph{descomposición} es lo que las salva: dividir el problema en pasos manejables convierte la parálisis en acción. Esa habilidad no es solo para programadores; \emph{es para vivir bien}. Aprenderla en 7° es regalo para tu yo adulto que enfrentará crisis.

\textbf{¿Qué problema actual mío se ha vuelto abrumador porque no lo he descompuesto?}
\end{softbox}
\linead{\linewidth}\\[1mm]
\linead{\linewidth}

\begin{softbox}{milcMostaza}{milcGris}{Estoicismo · Marco Aurelio · cuidado interior}
\textit{El imperio se gobierna provincia por provincia. La vida se vive día por día. Los problemas se resuelven sub-problema por sub-problema.}

Marco Aurelio gobernaba el imperio romano (millones de personas, miles de problemas simultáneos). Su técnica era \textbf{dividir}: cada problema en sub-problemas, cada sub-problema en tareas, cada tarea en un día. \emph{Quien quiere resolver todo a la vez, no resuelve nada. Quien divide, avanza}. Esa disciplina romana sobrevive en lo que hoy llamas pensamiento computacional.

\textbf{¿Estoy intentando resolver problemas grandes de un solo golpe? ¿Cómo cambiaría si los dividiera?}
\end{softbox}
\linead{\linewidth}\\[1mm]
\linead{\linewidth}

\begin{softbox}{milcTurquesa}{milcGris}{Floridi · lente de la infoesfera}
\textit{Quien aprende a pensar computacionalmente entiende el mundo. Quien no, sobrevive en el mundo de los demás.}

Las 4 técnicas son herramientas \textbf{para entender cualquier problema complejo}, no solo para programar. \emph{Los líderes que respetas} (médicos, ingenieros, profesores, emprendedores) aplican estas 4 técnicas a sus retos diarios, aunque no las llamen así. Aprenderlas explícitamente te coloca al lado de ellos en términos de cómo procesas el mundo.

\textbf{¿Estoy desarrollando una forma adulta de pensar problemas, o me quedo en reacciones emocionales?}
\end{softbox}
\linead{\linewidth}\\[1mm]
\linead{\linewidth}

\begin{infoband}{milcNegro}
\textbf{Nota para el docente.} Las tres formulaciones del triángulo son la síntesis del autor de la guía, no citas textuales de los pensadores. Si vas a citarlos en clase, remite a la obra original.
\end{infoband}

\milcestacion{milcVino}{Mi autoevaluación · marca una casilla por fila}

{\small
\setlength{\extrarowheight}{2.2pt}
\begin{tabularx}{\textwidth}{>{\RaggedRight\arraybackslash\bfseries}p{3.0cm}YYY}
\rowcolor{milcVino}\color{white}\bfseries Criterio & \color{white}\bfseries Lo logré & \color{white}\bfseries En proceso & \color{white}\bfseries Necesito apoyo\\[.6mm]
Comprensión del tema & \milcopcion{Explico las ideas con mis palabras.} & \milcopcion{Reconozco las ideas, falta precisión.} & \milcopcion{Necesito repasar lo básico.}\\[.8mm]
\rowcolor{milcGris}Pensamiento computacional & \milcopcion{Usé los pilares al armar mi producto.} & \milcopcion{Usé algunos pilares.} & \milcopcion{Necesito releer la estación 2.}\\[.8mm]
Aplicación y producto & \milcopcion{Mi producto cumple los criterios.} & \milcopcion{Está armado, con ajustes pendientes.} & \milcopcion{Aún está en construcción.}\\[.8mm]
\rowcolor{milcGris}Reflexión 5 dimensiones & \milcopcion{Respondí cada una con honestidad.} & \milcopcion{Respondí algunas con detalle.} & \milcopcion{Mis respuestas son muy generales.}\\[.8mm]
Triángulo de pensamiento & \milcopcion{Conecté las 3 voces con mi trabajo.} & \milcopcion{Conecté al menos una.} & \milcopcion{No las relaciono con mi proceso.}\\[.8mm]
\end{tabularx}}

\vspace{3mm}
\textbf{Hoy entendí que:}\\[1mm]
\linead{\linewidth}\\[1mm]
\linead{\linewidth}

\textbf{Mi compromiso para la próxima sesión es:}\\[1mm]
\linead{\linewidth}\\[1mm]
\linead{\linewidth}

\begin{softbox}{milcMostaza}{milcGris}{Cita de despedida · Marco Aurelio}
\textit{``El obstáculo en el camino se vuelve el camino. Lo que estorba al camino se vuelve camino.''}\\[1mm]
Cada error de hoy, bien observado, se convierte en pieza del camino que pisarás mañana.
\end{softbox}

\small
\begin{tabularx}{\textwidth}{>{\bfseries}p{3.6cm}Y}
\rowcolor{milcGradeDark!12}Nombre completo: & \linead{10cm}\\
Fecha: & \linead{4cm} \quad Curso/grupo: \linead{4cm}\\
Mi firma: & \linead{9cm}\\
\end{tabularx}
\normalsize

\begin{infoband}{milcGradeDark}
\textbf{Para la próxima sesión.} Esta semana, aplica las 4 técnicas a un problema más (puede ser un trabajo de otra materia o algo de tu vida). La próxima clase aprendes a \textbf{dibujar algoritmos}: los diagramas de flujo con símbolos estándar.
\end{infoband}

\end{document}
